% ===================================================================
% DRAFT -- new section for su7_hierarchy.tex, to be spliced after \section{What this predicts}
% (sec:pred) and before \section{A second anchor} (sec:anchor).
%
% Everything below is derived in collider_dictionary.py; its archived run is
% outputs/collider_dictionary.txt and every control there passes.
% ===================================================================
\section{Which state a search bounds}\label{sec:collider}

A bound on $1/R_5$ is a bound on a compactification scale, not on the mass of any one particle. A
collider does not measure compactification scales; it measures resonances, each with a production
rate and a width. Getting from one to the other looks like it needs a phenomenological study, and
it does not: it is fixed by the three parity matrices of \cite{KM25} and the mode expansions that
follow from them, all of which are already published. This section reads them.

\paragraph{The parity that separates colour.} Of the three $\ZZ_2$ matrices, the one acting in
$x_6$ is
\begin{equation}\label{eq:P6}
  P_6=\mathrm{diag}(+,+,+,-,-,-,-),
\end{equation}
\cite{KM25}'s (11), which splits the seven-dimensional index into the colour triplet and
everything else. A gauge component along $E_{ij}$ carries the \emph{product} $P_6^{(i)}P_6^{(j)}$,
so \eqref{eq:P6} has an immediate consequence that is arithmetic rather than dynamical:
\begin{equation}\label{eq:colourparity}
  P_6(i,j)=-1
  \quad\Longleftrightarrow\quad
  \text{exactly one of } i,j \text{ is a colour index,}
\end{equation}
that is, precisely when the generator changes colour. There are $24$ such components, and there is
no colour-changing component with $P_6=+1$.

\begin{theorem}[colour lives on the small circle]\label{thm:colourparity}
Every colour-changing gauge boson of \cite{KM25} has $P_6=-1$. By its (21)--(24) each such
component expands in $\sin(m x_6/R_6)$ with $m\ge1$. Hence its lightest level is at least $1/R_6$,
and its wavefunction \emph{vanishes} at $x_6=0$.
\end{theorem}
\begin{proof}
\eqref{eq:colourparity} is the statement that $P_6^{(i)}P_6^{(j)}=-1$ iff one index lies in the
colour block, which is immediate from \eqref{eq:P6}. The four $P_6$-odd expansions of \cite{KM25},
its (21)--(24), all carry $\sin(m x_6/R_6)$ and all start at $m=1$; the sine vanishes at the
origin and the level is bounded below by the $m=1$ term.
\end{proof}

\noindent This matters because \cite{KM25} p.~3 fixes $R_5\gg R_6$ and calls $1/R_6$ ``an
arbitrary parameter, which is the Planck scale for instance''. Theorem~\ref{thm:colourparity} then
says the coloured exotics are not merely heavy but unobservable in principle at a hadron collider,
and it says so without any assumption beyond the orbifold.

\paragraph{So the brane quarks see exactly two towers.} \cite{KM25} p.~22 introduces ``quarks
$q_L,u_R,d_R$ as the 4D fermions on the brane at the origin in a compactified space $(0,0)$''. A
4D field there couples only to components whose mode function is non-zero at that point --- that
is, to components whose $x_5$ \emph{and} $x_6$ factors are both cosines. Of the eight parity
classes only two qualify, and both are colour-diagonal:

\begin{center}\small
\rowcolors{2}{white}{softgrey}
\begin{tabular}{llllll}
\toprule
\rowcolor{hdrblue!10} parity & eq. & $x_5$ mode & $x_6$ mode & at $(0,0)$ & lightest level\\
\midrule
$(+,+,+)$ & (17) & $\cos(nx_5/R_5)$ & $\cos(mx_6/R_6)$ & \textbf{non-zero} & $0$, then $1/R_5$\\
$(+,+,-)$ & (18) & $\cos((n{+}\tfrac12)x_5/R_5)$ & $\cos(mx_6/R_6)$ & \textbf{non-zero} & $1/(2R_5)$\\
$(+,-,+)$ & (19) & $\sin((n{+}\tfrac12)x_5/R_5)$ & $\cos(mx_6/R_6)$ & zero & $1/(2R_5)$\\
$(+,-,-)$ & (20) & $\sin(nx_5/R_5)$ & $\cos(mx_6/R_6)$ & zero & $1/R_5$\\
$(-,\cdot,\cdot)$ & (21)--(24) & --- & $\sin(mx_6/R_6)$ & zero & $\ge 1/R_6$\\
\bottomrule
\end{tabular}
\end{center}

\noindent The $(+,+,-)$ class is exactly the two positions $(6,7)$ and $(7,6)$ --- a colour
singlet. So the half-integer tower, whose lightest state sits at $1/(2R_5)$ and which
\S\ref{sec:pred} correctly flagged as the lowest level in the model, is \emph{not} something a
dijet search can see; and it is not a clean prediction either, since $\langle W\rangle$ of
\cite{KM25}'s (77) acts in the $(5,7)$ block and shifts it with $\alpha$. The one coloured object
below the Planck scale is the gluon tower.

\paragraph{And that answers a question \cite{KM25} left open.} Their conclusions read: ``Since the
mass of $X$, $Y$ gauge bosons is small in our model, this causes rapid proton decay \emph{if they
can couple to quarks and leptons}'', and they leave the question to future work. By
Theorem~\ref{thm:colourparity} they cannot: the $X$ and $Y$ bosons are precisely the
colour-changing components, and their wavefunction is zero at the fixed point where the same paper
puts the quarks. The gauge-exchange channel is shut by the parity assignment itself.

\noindent\textbf{The scope of that is narrow and worth stating.} It is a tree-level statement about
gauge exchange, for the quarks localised as in \cite{KM25} p.~22. It says nothing about the brane
Dirac masses that same paper introduces on p.~16 to lift its unwanted zero modes, and it is
\emph{not} the escape of Part~VI \cite{PartVI}, which is a family-dependent charge hosted by a
donated $\rep{84}^{(+,+)}$ --- a different mechanism, addressing a different operator. In
particular \eqref{eq:escapeceiling} is untouched by it.

\paragraph{The coloron, with no free parameters.} Three facts fix it. The mass: \cite{KM25} p.~16
states that the $\su3_C$ generators ``are commutable with $\langle W\rangle$ irrespective of the
VEV $\alpha_{\min}$'', so the colour tower takes no Wilson-line shift and sits at $n/R_5$ exactly.
The coupling: their (17) carries $\sqrt2$ in front of every $(n,0)$ mode where the zero mode
carries $1$, and $\cos(0)=1$, so at the brane the ratio is $\sqrt2$ --- and it is
flavour-universal, because $q_L$, $u_R$ and $d_R$ sit at the same point and see the same mode
function. The width then follows from the standard colour-octet result
$\Gamma(C\to q\bar q)=n_f\,\alpha_s\cot^2\theta\,M/6$ with $\cot^2\theta=2$ and $n_f=6$:

\begin{keyeq}
\begin{equation}\label{eq:coloron}
  M_n=\frac{n}{R_5},
  \qquad
  \frac{g_{G^{(n)}q\bar q}}{g_{gq\bar q}}=\sqrt2,
  \qquad
  \frac{\Gamma}{M}=2\alpha_s(M)\simeq0.16 .
\end{equation}
\end{keyeq}

\noindent Carried across the rows of \S\ref{sec:pred}:

\begin{center}\small
\rowcolors{2}{white}{softgrey}
\begin{tabular}{lrrrr}
\toprule
\rowcolor{hdrblue!10} & $1/R_5$ & $\alpha_s$ & $\Gamma/M$ & $\Gamma$\\
\midrule
\rowcolor{amber!14}ceiling at the measured $m_h$ & $9.09$~TeV & $0.0735$ & $0.147$ & $1337$~GeV\\
ceiling once the vacuum must be the true one & $9.22$~TeV & $0.0734$ & $0.147$ & $1354$~GeV\\
ceiling over \emph{all} contents & $10.03$~TeV & $0.0729$ & $0.146$ & $1463$~GeV\\
contents that can afford the escape & $3.97$~TeV & $0.0789$ & $0.158$ & $626$~GeV\\
\bottomrule
\end{tabular}
\end{center}

\paragraph{Which changes what may be quoted against it.} A width of one sixth of the mass is not a
narrow resonance, so the narrow benchmark of a dijet search does not bound this object: CMS's
$6.6$~TeV for a narrow axigluon or coloron at $137\,\mathrm{fb}^{-1}$ \cite{CMSdijet} is a limit on
a different particle. What does apply is the model-independent limit on $\sigma\,B\,A$ that the
same analysis publishes as a function of mass \emph{and} width, read at $\Gamma/M\simeq0.16$. The
previous version of this section compared \eqref{eq:escapeceiling} to a generic Kaluza-Klein gluon
exclusion; \eqref{eq:coloron} is what makes that comparison a statement about this model rather
than about a benchmark, and it is the reason the comparison is worth making at all.

\paragraph{A correction to how the reach was stated.} An earlier version of this paper wrote that a
machine reaching $10$--$20$~TeV settles the class. That conflates the energy of a collider with its
reach in resonance mass, and the two are not the same: producing a $9$~TeV resonance at
$\sqrt s=13$~TeV requires partons at values of $x$ where the parton densities are negligible. The
correct statement is about mass reach in the relevant channel, and it makes the escape branch far
more interesting than the ceiling --- $3.97$~TeV sits where Run~2 has real sensitivity, while
$9.09$~TeV does not.

\paragraph{The tower also acts below the first resonance.} Since every $(n,0)$ mode has the same
$\sqrt2$ at $x_5=0$, integrating out the whole tower gives a colour-octet four-quark operator with
a coefficient that is a pure number:
\begin{equation}\label{eq:tower}
  \sum_{n\ge1}\frac{g_n^2}{2M_n^2}
  \;=\;g_s^2R_5^2\sum_{n\ge1}\frac{1}{n^2}
  \;=\;\frac{\pi^2}{6}\,g_s^2R_5^2 ,
  \qquad\text{i.e.}\qquad
  \Lambda_8=\frac{\sqrt3}{\pi\sqrt{\alpha_s}}\,\frac{1}{R_5},
\end{equation}
which is $7.8$~TeV on the escape branch and $18.5$~TeV at the measured-$m_h$ ceiling. The $m\ge1$
modes do not spoil the sum: they cost $1/R_6$, and their share is suppressed by $R_6/R_5$, which
\cite{KM25} p.~3 takes to be tiny --- so the logarithmic divergence one would fear from a
codimension-two brane never switches on at the scales in play. This is a collective prediction of
the tower and does not require resolving any resonance; it cannot, however, be compared directly
with published compositeness scales, whose operator has a different colour and chiral structure.

\paragraph{And one number the model does not supply.} \cite{KM25} p.~16 lifts its unwanted zero
modes by introducing ``4D fermion conjugate to the representations and charges of the massless
fermions and their Dirac mass terms on the fixed points'', and does not fix those masses. If any
coloured exotic sits below $M/2$ then $G^{(1)}\to\Psi\bar\Psi$ opens, the width grows and
$\mathrm{BR}(G^{(1)}\to jj)$ falls by an amount nothing here determines. So
\begin{equation}\label{eq:brhole}
  \mathrm{BR}(G^{(1)}\to jj)=1 \quad\text{if every coloured exotic lies above } M/2,
\end{equation}
and less otherwise. That single unknown is what stands between the arithmetic ceiling and a
model-specific exclusion, and a limit quoted without it would be model-independent in appearance
only. It is stated here rather than absorbed.

\bridge{The dictionary is derived, the one free parameter is named, and the comparison is now a
statement about this model --- which leaves only the arithmetic instrument itself to be checked
against someone else's table.}
